\documentclass[12pt,a4paper,titlepage]{article}
\usepackage{parskip}  %gestion des espaces entre paragraphes
\usepackage{enumitem}
\def \verseSpace{\vskip 0.45em}
\def \parSpace{\vskip 1.2em}
\def \decalage{\hspace*{0.5cm}~}
\newlength{\indentunit}
\setlength{\indentunit}{0.5cm}
\newlength{\curindent}
\setlength{\curindent}{0pt}
\newcommand{\indentlvl}[1]{\setlength{\curindent}{#1\indentunit}\hspace*{\curindent}}
\newcommand{\brk}{\\
\hspace*{\curindent}} % remplace \\ dans les paragraphes indentés
\def \pscolor{red} %couleur de la numérotation des versets
\usepackage{titlesec} %gestion des formats de titre
\usepackage{xcolor}
\usepackage{graphicx}
\usepackage{geometry}
\usepackage{fancyhdr}
%\usepackage[french]{babel}

\usepackage{fontspec}
\usepackage{polyglossia}
\setdefaultlanguage{french}
\setotherlanguages{hebrew}
\newfontfamily\hebrewfont[Script=Greek]{SBL Greek}
%\newfontfamily\libertinefont{Linux Libertine O}
\usepackage{libertine} %police utilisée

\geometry{
  twoside=false,
  left=10mm,
  right=10mm,
  top=15mm,
  bottom=15mm
}
  
  \titleformat{\section}
  {\Large\bfseries\scshape}{\thesection}{5em}{}



%%%%% variables pour la numérotation des versets %%%%%%
\def \pslabelsep{0.3em} %distance entre le numéro de verset et le texte
\def \psleftmargin{1.5em} %marge à gauche lors de la numérotation des versets de psaumes

\begin{document}

\pagestyle{fancy}
\fancyhf{} % Efface les en-têtes et pieds de page par défaut
\fancyfoot[R]{\footnotesize{}} % Affiche le numéro de page centré dans le pied de page
\renewcommand{\headrulewidth}{0pt} % Supprime la ligne horizontale en haut de chaque page

\noindent
\section*{}
\subsection*{}
\vspace{-0.62mm}

\begin{enumerate}[leftmargin=\psleftmargin, labelsep=\pslabelsep, label=\arabic*, align=right, font=\fontspec{Linux Libertine}\color{\pscolor}\small\textsuperscript, parsep=0em, itemsep=0em, topsep=0em]
  \item Παῦλος ἀπόστολος \\
        \indentlvl{1}οὐκ ἀπ’ ἀνθρώπων\brk
        οὐδὲ δι’ ἀνθρώπου \brk
        ἀλλὰ \\
        \indentlvl{2}διὰ Ἰησοῦ Χριστοῦ \brk
        καὶ θεοῦ πατρὸς τοῦ \\
        \indentlvl{3}ἐγείραντος τὸν ἐκ νεκρῶν,
  \item καὶ οἱ σὺν ἐμοὶ πάντες ἀδελφοὶ \\
        ταῖς ἐκκλησίαις τῆς Γαλατίας,
  \item χάρις ὑμῖν καὶ εἰρήνη \\
        \indentlvl{1}ἀπὸ θεοῦ πατρὸς ἡμῶν \brk
        καὶ κυρίου Ἰησοῦ Χριστοῦ
  \item \indentlvl{2}τοῦ δόντος ἑαυτὸν ὑπὲρ τῶν ἁμαρτιῶν ἡμῶν,\\
        \indentlvl{3}ὅπως ἐξέληται ἡμᾶς ἐκ τοῦ αἰῶνος τοῦ ἐνεστῶτος πονηροῦ \brk
        κατὰ τὸ θέλημα τοῦ θεοῦ καὶ πατρὸς ἡμῶν,
  \item  \indentlvl{4}ᾧ ἡ δόξα εἰς τοὺς αἰῶνας τῶν αἰώνων, ἀμήν. \parSpace

  \item Θαυμάζω ὅτι οὕτως ταχέως μετατίθεσθε \\
        \indentlvl{1}ἀπὸ τοῦ καλέσαντος ὑμᾶς ἐν χάριτι [Χριστοῦ]\\
        \indentlvl{2}εἰς ἕτερον εὐαγγέλιον,
  \item \indentlvl{3}ὃ οὐκ ἔστιν ἄλλο, \\
        \indentlvl{4}εἰ μή τινές εἰσιν οἱ ταράσσοντες ὑμᾶς \\
        \indentlvl{4}καὶ θέλοντες μεταστρέψαι \\
        \indentlvl{5}τὸ εὐαγγέλιον τοῦ Χριστοῦ.
  \item ἀλλὰ καὶ ἐὰν \\
        \indentlvl{1}ἡμεῖς ἢ ἄγγελος ἐξ οὐρανοῦ\\
        εὐαγγελίζηται [ὑμῖν] \\
        \indentlvl{1}παρ’ ὃ εὐηγγελισάμεθα ὑμῖν, \\
        ἀνάθεμα ἔστω.
  \item ὡς προειρήκαμεν καὶ ἄρτι πάλιν λέγω·\\
        \indentlvl{1} εἴ τις ὑμᾶς εὐαγγελίζεται παρ’ ὃ παρελάβετε,\\
        ἀνάθεμα ἔστω.
  \item Ἄρτι γὰρ ἀνθρώπους πείθω ἢ τὸν θεόν;\\
        ἢ ζητῶ ἀνθρώποις ἀρέσκειν;\\
        εἰ ἔτι ἀνθρώποις ἤρεσκον,\\
        \indentlvl{1} Χριστοῦ δοῦλος οὐκ ἂν ἤμην.
        \parSpace

  \item Γνωρίζω γὰρ ὑμῖν, ἀδελφοί,\\
        \indentlvl{1}τὸ εὐαγγέλιον τὸ εὐαγγελισθὲν ὑπ’ ἐμοῦ\brk
        ὅτι οὐκ ἔστιν κατὰ ἄνθρωπον·
  \item οὐδὲ γὰρ ἐγὼ παρὰ ἀνθρώπου παρέλαβον αὐτὸ\\
        \indentlvl{1} οὔτε ἐδιδάχθην ἀλλὰ\\
        \indentlvl{2}δι’ ἀποκαλύψεως Ἰησοῦ Χριστοῦ.
  \item Ἠκούσατε γὰρ τὴν ἐμὴν ἀναστροφήν ποτε ἐν τῷ Ἰουδαϊσμῷ,\\
        \indentlvl{1}ὅτι καθ’ ὑπερβολὴν ἐδίωκον τὴν ἐκκλησίαν τοῦ θεοῦ \brk
        καὶ ἐπόρθουν αὐτήν,
  \item \indentlvl{1}καὶ προέκοπτον ἐν τῷ Ἰουδαϊσμῷ \\
        \indentlvl{2}ὑπὲρ πολλοὺς συνηλικιώτας ἐν τῷ γένει μου, \brk
        περισσοτέρως ζηλωτὴς ὑπάρχων τῶν πατρικῶν μου παραδόσεων.
  \item Ὅτε δὲ εὐδόκησεν [ὁ θεὸς]  \\
        \indentlvl{1}ὁ ἀφορίσας με ἐκ κοιλίας μητρός μου \brk
        καὶ καλέσας διὰ τῆς χάριτος αὐτοῦ 
  \item \indentlvl{2}ἀποκαλύψαι τὸν υἱὸν αὐτοῦ ἐν ἐμοί,\brk
  \indentlvl{2}ἵνα εὐαγγελίζωμαι αὐτὸν ἐν τοῖς ἔθνεσιν, \\
  εὐθέως οὐ προσανεθέμην σαρκὶ καὶ αἵματι
  \item \indentlvl{1}οὐδὲ ἀνῆλθον εἰς Ἱεροσόλυμα πρὸς τοὺς πρὸ ἐμοῦ ἀποστόλους, \\
  ἀλλ’ ἀπῆλθον εἰς Ἀραβίαν \\
  \indentlvl{1}καὶ πάλιν ὑπέστρεψα εἰς Δαμασκόν.
  \item Ἔπειτα μετὰ ἔτη τρία ἀνῆλθον εἰς Ἱεροσόλυμα\\
  \indentlvl{1}ἱστορῆσαι Κηφᾶν \brk
  καὶ ἐπέμεινα πρὸς αὐτὸν ἡμέρας δεκαπέντε,
  \item ἕτερον δὲ τῶν ἀποστόλων οὐκ εἶδον\\
  \indentlvl{1}εἰ μὴ Ἰάκωβον τὸν ἀδελφὸν τοῦ κυρίου.
  \item ἃ δὲ γράφω ὑμῖν, ἰδοὺ ἐνώπιον τοῦ θεοῦ\\
  \indentlvl{1}ὅτι οὐ ψεύδομαι.
  \item Ἔπειτα ἦλθον εἰς τὰ κλίματα τῆς Συρίας καὶ τῆς Κιλικίας·
  \item ἤμην δὲ ἀγνοούμενος τῷ προσώπῳ\\
  \indentlvl{1}ταῖς ἐκκλησίαις τῆς Ἰουδαίας ταῖς ἐν Χριστῷ.
  \item μόνον δὲ ἀκούοντες\\
  \indentlvl{1}ἦσαν ὅτι ὁ διώκων ἡμᾶς ποτε\brk
  νῦν εὐαγγελίζεται τὴν πίστιν ἥν ποτε ἐπόρθει,
  \item καὶ ἐδόξαζον ἐν ἐμοὶ τὸν θεόν.
\end{enumerate}
\end{document}